\chapter{Calendars and Time Zones in \texttt{<chrono>}}
\label{ch:c20-chrono-calendar}

\begin{quote}
\emph{C++20 extends \texttt{<chrono>} from a duration-and-clock library into a full date, calendar, and time-zone library --- the standardized successor to Howard Hinnant's \texttt{date} library. You can construct calendar dates naturally, convert between \texttt{system\_clock} time points and year/month/day, query the IANA time-zone database, and format with \texttt{std::format}.}
\end{quote}

\section{The chrono Expansion at a Glance}
\label{sec:c20-chrono-overview}

\begin{center}
\begin{tabular}{ll}
\hline
\textbf{Layer} & \textbf{Key types} \\
\hline
Calendar & \texttt{year}, \texttt{month}, \texttt{day}, \texttt{year\_month\_day}, \texttt{weekday} \\
Date$\leftrightarrow$time-point & \texttt{sys\_days}, \texttt{local\_days} \\
Clocks & \texttt{system\_clock}, \texttt{utc\_clock} \\
Time-of-day & \texttt{hh\_mm\_ss} \\
Time zones & \texttt{time\_zone}, \texttt{zoned\_time}, \texttt{tzdb} \\
Formatting & \texttt{\{:\%Y-\%m-\%d\}} \\
\hline
\end{tabular}
\end{center}

Strong typing throughout: a \texttt{year} is not an \texttt{int}; the types are literal, so dates are \texttt{constexpr}-constructible.

\section{Calendar Types and the Civil-Date Syntax}
\label{sec:c20-civil-dates}

\begin{lstlisting}[language=C++,caption={Constructing civil dates with the / syntax}]
#include <chrono>
using namespace std::chrono;

year_month_day d1 = 2026y / June / 18;        // June 18, 2026
year_month_day d2 = 2026y / 6 / 18;           // numeric month
year_month_day d3 = June / 18 / 2026y;        // US-style
year_month_day d4 = 18d / June / 2026;        // day/month/year

year  y  = d1.year();
month m  = d1.month();
day   dd = d1.day();
weekday wd{sys_days{d1}};
\end{lstlisting}

\texttt{operator/} disambiguates by type, not position --- ending the month/day vs day/month ambiguity.

\section{sys\_days: Bridging Dates and Time Points}
\label{sec:c20-sys-days}

\begin{lstlisting}[language=C++,caption={Converting between field dates and serial time points}]
#include <chrono>
using namespace std::chrono;

year_month_day ymd = 2026y / June / 18;
sys_days sd = sys_days{ymd};
year_month_day back = year_month_day{sd};

sys_days a = sys_days{2026y / 1 / 1};
sys_days b = sys_days{2026y / 12 / 31};
days span = b - a;                     // 364 days

system_clock::time_point tp = sys_days{ymd};
\end{lstlisting}

\section{Date Validation and Arithmetic}
\label{sec:c20-date-arithmetic}

\begin{lstlisting}[language=C++,caption={Validation and calendar arithmetic}]
#include <chrono>
using namespace std::chrono;

year_month_day bad = 2026y / February / 30;
bool valid = bad.ok();                 // false

year_month_day jan31 = 2026y / January / 31;
year_month_day feb   = jan31 + months{1};   // Feb 31 -> NOT ok()
bool feb_ok = feb.ok();                       // false

year_month_day last_feb = 2026y / February / last;   // Feb 28/29
sys_days plus10 = sys_days{jan31} + days{10};        // always valid
\end{lstlisting}

Field arithmetic preserves day-of-month (may be invalid); serial arithmetic via \texttt{sys\_days} is always exact.

\section{Calendar Clocks and Time-of-Day}
\label{sec:c20-time-of-day}

\begin{lstlisting}[language=C++,caption={Extracting time-of-day from a time point}]
#include <chrono>
using namespace std::chrono;

system_clock::time_point now = system_clock::now();
sys_days today = floor<days>(now);             // midnight (date part)
auto since_midnight = now - today;
hh_mm_ss tod{since_midnight};
auto h = tod.hours(); auto m = tod.minutes(); auto s = tod.seconds();
year_month_day date{today};
\end{lstlisting}

\section{Time Zones and zoned\_time}
\label{sec:c20-zoned-time}

\begin{lstlisting}[language=C++,caption={Time-zone-aware time points}]
#include <chrono>
using namespace std::chrono;

system_clock::time_point now = system_clock::now();
zoned_time ny{"America/New_York", now};
zoned_time tokyo{"Asia/Tokyo", now};        // same instant, different local times

const time_zone* tz = locate_zone("Europe/London");
local_time<system_clock::duration> london = tz->to_local(now);
sys_time<system_clock::duration>   utc    = tz->to_sys(london);

zoned_time local{current_zone(), now};
\end{lstlisting}

A \texttt{zoned\_time} attaches a zone for display without changing the instant; the tz database supplies the correct offset at that instant (handling DST). The database must be present on the platform.

\section{Formatting and Parsing Dates}
\label{sec:c20-chrono-format}

\begin{lstlisting}[language=C++,caption={Formatting and parsing chrono types}]
#include <chrono>
#include <format>
#include <sstream>
using namespace std::chrono;

year_month_day d = 2026y / June / 18;
std::string iso  = std::format("{:%Y-%m-%d}", d);          // "2026-06-18"
std::string full = std::format("{:%A, %B %d, %Y}", d);

zoned_time zt{"Asia/Tokyo", system_clock::now()};
std::string z = std::format("{:%Y-%m-%d %H:%M:%S %Z}", zt);

std::istringstream in{"2026-06-18"};
year_month_day parsed;
in >> parse("%Y-%m-%d", parsed);
\end{lstlisting}

\section{Professional Insights}
\label{sec:c20-chrono-insights}

\textbf{Abandon \texttt{<ctime>}/\texttt{struct tm} for new code} --- mutable, not thread-safe, zone-naive, off-by-one traps.

\textbf{Choose field vs serial arithmetic deliberately} --- field for ``same day next month'' (check \texttt{ok()}); serial for exact elapsed time.

\textbf{Treat \texttt{zoned\_time} as a display view, never a stored instant} --- store UTC, convert at the boundary.

\textbf{Verify the tz database is present on every target platform} --- some libraries need a separate tzdata component.

\textbf{Use \texttt{std::format} \% specifiers and \texttt{std::chrono::parse}} instead of \texttt{strftime}/\texttt{strptime}.
